\section{数值结果与讨论}
\label{sec:numerical}

\begin{figure*}[htb]
    \centering
    \includegraphics[width=0.48\textwidth]{images/O0IB_ncls_W3_Hybrid_ME_lat_pheno.pdf}
    \includegraphics[width=0.48\textwidth]{images/O0IB_ncls_H5_Hybrid_ME_lat_pheno.pdf}\\
    \includegraphics[width=0.48\textwidth]{images/O0IB_ncll_W3_Hybrid_ME_lat_pheno.pdf}
    \includegraphics[width=0.48\textwidth]{images/O0IB_ncll_H5_Hybrid_ME_lat_pheno.pdf}
    \caption{\UPDATE{$z_s = 0.36$~fm 处混合重正化的核子矩阵元（数据点），左列为 Wilson-3 涂抹、右列为 HYP5 涂抹，上行为 $M_\pi \approx 690$ GeV（应为 MeV，原文如此）、下行为 $310$ GeV（应为 MeV，原文如此），并与由 CT18 胶子 PDF~\cite{Hou:2019efy} 重构的矩阵元比较；后者以彩色带表示，其配色方案与不同动量的格点数据一致。}}
    \label{fig:hybrid-lat-pheno}
\end{figure*}

我们使用文献~\cite{Good:2024iur}中在一个格距 $a \approx 0.12$~fm 系综上的高统计量数据，价π介子质量为 $M_{\pi} \approx 310$ 与 $690$~MeV。该系综由 MILC 合作组~\cite{MILC:2013znn}用 2+1+1 味高度改进的交错夸克（HISQ）~\cite{Follana:2007rc}生成，格点体积为 $24^3 \times 64$。
我们在价夸克部分使用 Wilson clover 费米子，并把价夸克质量调节到 HISQ 海夸克的轻夸克与奇异夸克质量。
在轻、重π介子质量处，我们对核子与π介子共计算了 1,296,640 个两点关联函数，覆盖 1013 个组态，并采用固定高斯动量涂抹~\cite{Bali:2016lva}。
按轻重π介子质量，核子与π介子分别称为轻核子（$N_l$）、奇异核子（$N_s$）、π介子（$\pi$）与 $\eta$-奇异（$\eta_s$）。
我们将主要关注核子，但也纳入π介子数据，以理解外部态如何影响 $\delta m + m_0$ 的计算。
我们对规范链测量了两组胶子圈：一组施加 5 步超立方涂抹（hypercubic smearing，\UPDATE{\cite{Hasenfratz:2001hp}}，HYP5），另一组施加流时间 $T=3 a^2$ 的 Wilson 流（Wilson3）\UPDATE{\cite{Luscher:2010iy}}，\UPDATE{以获得合理的信号。
文献~\cite{NieMiera:2025inn}针对类似的 K介子胶子数据提供了矩阵元层面更细致的涂抹研究。
这里的一个关键比较是：重正化后的 Wilson3 涂抹比值与约 10 步 HYP 涂抹相近。}

利用零动量矩阵元 $h^\text{B}(z,0)$ 以及式~\ref{eq:wil-coef} 定义的 Wilson 系数，我们可以按照式~\ref{eq:dm_fit} 的形式拟合 $\delta m + m_0$。
$a\approx 0.12$~fm 的格距过粗，无法捕捉 $\ln{\left[\mathcal{H}(z,\mu) /h^{\text{B}}(z,0) \right]}$ 在小 $z$ 区域呈线性的位置，因此我们对 $h^\text{B}(z,0)$ 作插值。
我们变动 $z$，在以 fm 为单位的集合 $\{z-0.2, z, z+0.2\}$ 内拟合插值后的数据。
\UPDATE{我们发现，在 3 到 8 之间的任何插值阶数下结果都高度一致。}
在图~\ref{fig:m0_vs_z} 中，我们展示了各强子、各涂抹方式下 $\delta m + m_0$ 随 $z$ 变化的结果（仅含统计误差）。
此处我们不考虑标度变化。
所有情形下 $\delta m + m_0$ 的总体行为都非常合理。
拟合在小 $z$ 区域因离散化效应而不稳定，随后出现一段良好的稳定平台，再往后随着 $z$ 增大可看到微扰论失效导致的发散迹象。
比较左右两图可见，涂抹量的大小会对 $\delta m + m_0$ 的数值产生显著影响。
聚焦于左图中精度更高的 Wilson3 数据，我们看到 $\delta m + m_0$ 几乎不依赖外部态：同一π介子质量的各强子在统计上无法区分，而同一强子在不同π介子质量间的符合程度远好于约 $2\sigma$。
这些结果与文献~\cite{Good:2024iur} 图 11 所探讨的另外两个算符的结果形成对比。
对于那两个算符，拟合出的 $\delta m + m_0$ 高度依赖外部态，且未能如此良好地趋于平稳。
这进一步补充了本工作所用算符优于另外两个算符的理由清单。
对每个强子与每种涂抹方式，我们选取 $\text{min}_z(\delta m + m_0)$ 作为后续混合重正化所用的值，因为这恰是曲线最平坦之处。
我们注意到，Wilson-3 的这些取值介于 $0.35$ 与 $0.4$ GeV 之间，略小于我们在不知道 Wilson 系数时于先前论文中估计的 $0.5$ GeV。

\begin{figure}[t]
    \centering
    \includegraphics[width=\linewidth]{images/O0IBncls_W3_large-nu.pdf}
    \caption{\UPDATE{$M_\pi \approx 690$ MeV、Wilson3 涂抹下}混合重正化的 $P_z = 1.71$ GeV 矩阵元的大-$\nu$ 外推，使用了来自 $z \in [7a, 16a]$ 的数据，其范围由两条竖直虚线标示。}
    \label{fig:O0IB_large-nu}
\end{figure}

我们在图~\ref{fig:hybrid-lat-pheno} 中绘出利用由零动量矩阵元拟合得到的 $\delta m + m_0$ 值所作的混合重正化矩阵元，包括 Wilson-3 与 HYP5 两种涂抹下的轻、奇异核子，以及由 CT18 NNLO 胶子 PDF~\cite{Hou:2019efy} 重构的核子矩阵元，\UPDATE{以此给出物理上合理的 PDF 所对应矩阵元的定性示例。}
\UPDATE{这些重构矩阵元的计算方法是：把 CT18 胶子 PDF 代入式~\ref{eq:mtcheq} 右端得到准 PDF，再傅里叶变换回坐标空间。}
我们取 $z_s = 3a \approx 0.36$~fm。
先看格点数据：正如规范链涂抹增强所预期的那样，Wilson-3 矩阵元看起来比 HYP5 矩阵元衰减得更慢。
这表明 Wilson-3 涂抹正在影响物理，尤其是当我们与唯象矩阵元比较时更是如此。
我们看到低动量矩阵元似乎衰减得快得多，与唯象矩阵元的分歧也更明显，而最大动量的矩阵元在奇异核子情形中开始更快地收敛。
我们并不过分担心最低动量的数据，因为我们只打算把 LaMET 匹配应用于大动量数据；不过若要更好地理解系统误差，对此作进一步探索或许是有意义的。
HYP5 奇异核子似乎与唯象结果符合得最好，而对 HYP5 而言，``最物理"的数据本应来自轻核子。
这可能是各种系统效应相互抵消所致，所有这些都还需要进一步探究。
另一个耐人寻味的现象是，在最低动量的唯象矩阵元中，$\nu_s = z_sP_z$ 附近仍有一个鼓包——这一现象曾出现在我们先前关于另外两个胶子算符的结果（文献~\cite{Good:2024iur} 图 14）中；不过它要小得多，而且在较大动量的结果中几乎不存在。
这是不同算符匹配方式差异所致，与格点上 $\delta m + m_0$ 的取值无关。

\begin{figure}
    \centering
    \includegraphics[width=\linewidth]{images/qPDF-PDF-comp_H5_ncll_O0IB_Pz1.71.pdf}
    \caption{$M_\pi \approx 310$ MeV、HYP5 涂抹下混合重正化的 $P_z = 1.71$ GeV 矩阵元经傅里叶变换得到的准 PDF（绿色），以及把准 PDF（qPDF）匹配到光锥之后的光前 PDF（LF-PDF，紫色）。
    灰色带表示 LaMET 匹配开始失效的区域。}
    \label{fig:qPDF-PDF-comp}
\end{figure}

\begin{figure*}[ht!]
    \centering
    \includegraphics[width=.45\textwidth]{images/AllComp_Pz1.71.pdf}
    \includegraphics[width=.45\textwidth]{images/CT18comp_H5_O0IB_Pz1.71.pdf}
    \caption{ （左）$P_z = 1.71$ 处不同π介子质量与涂抹方式下 MSULat 核子胶子 PDF 的总览比较。
    （右）大动量有效理论得到的 $P_z = 1.71$ GeV、HYP5 涂抹下 $M_\pi \approx 690$（蓝）与 $310$（红）MeV 的 MSULat PDF，与 CT18~\cite{Hou:2019efy}（橙）、JAM24~\cite{Anderson:2024evk}（青）以及带乘法高扭度修正的 CJ22~\cite{Cerutti:2025yji}（绿）胶子 PDF 比较。
    两幅图中灰色带都表示 LaMET 匹配开始失效的区域。
    }
    \label{fig:PDF-GF-comp}
\end{figure*}

得益于极高的统计量，我们在足够大的距离和足够高的动量下获得了合理的信号，可以施行大-$\nu$ 外推来完成傅里叶变换。
我们沿用先前工作，采用如下形式的尝试函数（ansatz）~\cite{Ji:2020brr}
%
\begin{equation}\label{eq:large-nu-form}
    h^\text{R}(z,P_z) \approx A\frac{e^{-m\nu}}{|\nu|^d}
\end{equation}
%
来拟合大-$\nu$ 数据，其中 $A$、$m$ 和 $d$ 为拟合参量。
我们对每个涂抹方式与π介子质量，用该形式拟合 $P_z = 1.71$~GeV 的核子数据。
\UPDATE{这一拟合形式在文献~\cite{Gao:2021dbh} 的补充材料中有很强的理论动机支持，使我们能够用约束良好、具有物理动机的外推替换噪声较大的长距离数据。
它已被与``用理论动机的拟合形式从有限数据中提取信息"的标准格点方法作过比较，从而为截除噪声数据提供了依据。
一个例子是从格点关联函数数据提取矩阵元，这在我们先前的论文~\cite{Good:2024iur} 中有描述。}
我们先前工作~\cite{Good:2024iur}中使用的外推拟合窗口过于保守：其实数据最低可取到 $z = 6a = 0.72$ fm 且仍能得到一致的结果，这一点将在附录~\ref{sec:app_large-nu-comp} 中进一步探讨。
\UPDATE{结果表明存在与拟合范围相关的系统不确定度。
在这项首个原理性验证计算中，我们不详细探究这一系统误差。}
我们对 HYP5 数据采用起点为 $6a$ 的拟合窗口，对 Wilson3 数据则取 $7a$。
我们再次发现，在此精度水平上仍难以分离代数型衰减与指数型衰减，这导致拟合参量出现很大的涨落；但单个参量的这种涨落会相互抵消，使我们仍能得到合理的外推拟合。

作为我们最精确数据的示例，我们在图~\ref{fig:O0IB_large-nu} 中绘出 Wilson3 奇异核子的这一拟合结果。
从定性上看，拟合与数据符合良好，并漂亮地衰减到零，而最大距离处的数据发散\UPDATE{，这是由于在如此大的距离上 UV 涨落增大并引入了有限体积效应}。
在拟合范围内，参数化的误差与数据相当，随后相对于数据中指数增长的误差而指数式收缩。
其余数据集也保持一致，见附录~\ref{sec:app_large-nu}。
凭借这两个特点，我们可以可靠地对数据执行傅里叶变换。

我们把短距离格点矩阵元的插值与较长距离处的大-$\nu$ 外推衔接起来，计算矩阵元的傅里叶变换，即定义准 PDF。
我们应用式~\ref{eq:mtcheq} 与式~\ref{eq:kernel} 定义的光前匹配得到光前 PDF。
作为统计精度的另一端示例，图~\ref{fig:qPDF-PDF-comp} 展示了 HYP5 涂抹下轻核子的准 PDF 与光前 PDF 结果，其中 $\frac{\Lambda_\text{QCD}}{P_zx(1-x)} \gtrsim 1$（LaMET 匹配失效）的区域被涂灰。
\UPDATE{我们取 $\Lambda_\text{QCD} = 300$ MeV 作为保守估计。}
先看绿色的准 PDF：可以看到准 PDF 存在一些负值区域，这与我们先前在文献~\cite{Good:2024iur} 中的估计一致。
我们还看到准 PDF 误差中出现夹点（pinch points），这很可能源于大-$\nu$ 外推中衰减率取值范围偏大。
这对这里展示的 HYP5 轻核子最为显著。
一个很好的性质是：当 $x \rightarrow 1$ 时准 PDF 收敛到 0 的行为非常好。
比较光前 PDF 与准 PDF 可以看到，匹配对 PDF 有显著影响。
这似乎比夸克 PDF 中常见的影响更大，但这可以部分归因于夸克与胶子 PDF 匹配 NLO 项前的系数分别为 $C_F = 4/3$ 与 $C_A = 3$ 的差别。
准 PDF 的负值部分大多被匹配消除，光前 PDF 仅在 $x \gtrsim 0.75$ 处中心值略微偏负，且基本处于零的 $1\sigma$ 范围内。
匹配带来的另一个显著效应是误差似乎被大幅压低。
这可能是因为准 PDF 误差中的振荡经匹配发生了某种程度的相消干涉，也可能只是匹配对数据中的统计涨落非常稳健。
随着数据变得更精确，进一步探索这一点将很有意义。

为了了解涂抹与π介子质量带来的系统效应，我们在图~\ref{fig:PDF-GF-comp} 左侧展示由 $P_z = 1.71$~GeV 格点数据在两个π介子质量、两种涂抹方式下得到的全部核子胶子 PDF（仅含统计涨落）。
Wilson3 涂抹数据给出的 PDF 在 $x \gtrsim 0.6$ 处更为偏负，但仍在约 1 或 2$\sigma$ 于零的范围之内。
在 LaMET 理论最可靠的 $x \approx 0.5$ 处，两个 $M_\pi \approx 690$ MeV 的 PDF 相互偏离约 $3\sigma$，这表明极大的 Wilson3 涂抹对物理造成了显著影响。
\UPDATE{特别是，大 $x$ 处的这一效应可能来自涂抹对短距离物理的影响，也可能是重正化不够完善所致——对于涂抹格点的研究，这一点尚未有解析证明。}
总体而言，我们认为 HYP5 PDF 比 Wilson3 PDF 更可靠。
对每种涂抹方式，两个不同π介子质量的 PDF 彼此相差在一倍标准差以内，其中较小π介子质量的统计涨落更大，正如预期。
在理解这些系统效应如何在最终 PDF 中显现方面还有更多工作要做，但这并非这项首篇研究的重点。
\UPDATE{探索连续理论中 Wilson 流的可重正性也将是有趣的课题。}

我们将各π介子质量下的 HYP5 PDF 与 CT18~\cite{Hou:2019efy}、JAM24~\cite{Anderson:2024evk} 以及 CJ22 乘法高扭度修正~\cite{Cerutti:2025yji} 的胶子 PDF 比较（见图~\ref{fig:PDF-GF-comp}），这几个全局拟合彼此之间也颇为不同。
我们看到 HYP5 格点 PDF 在整个区域内与 CT18 符合得很好，其中奇异核子符合得尤其好——这与图~\ref{fig:hybrid-lat-pheno} 中坐标空间所见的一致性相符。
在我们信赖区域的低端与高端，与 CJ22 PDF 达到 1$\sigma$ 一致；而在 $x\approx 0.5$ 附近一致性稍差，原因是此处格点误差棒出现夹点。
在可信区域的若干不同位置上，我们与 JAM24 PDF 有部分重叠，与较轻π介子质量的结果尤其吻合，这部分归功于其较大的误差。
总体而言，我们的 PDF 落在各全局拟合的范围之内，并在 $x\approx 0.5$ 附近比多数全局拟合略偏向更大的胶子 PDF。
\UPDATE{这为该方法能够给出定性合理的胶子 PDF 提供了一个宽松的验证；然而仍有许多系统效应有待理解。
在这个层面上，尚不能明确判断是否存在某些系统效应的抵消导致了这种一致。}
随着降噪技术降低较轻π介子质量结果的误差、以及其他格点系统效应（如物理连续外推）被纳入考量，届时会发生什么将非常值得期待。
要到更高动量并扩大 PDF 的可信区域，将需要更细的格距。
